% GENERATED FILE. Edit canonical lessons, then run npm run generate:books.
% canonical-source-hash: fnv1a64:23818d511ddda8ec
% canonical-lessons: KA-C495-muddisu, KA-C495-muttidu, KA-C495-vatsalya, KA-C495-mamate, KA-C495-kalaji

\chapter{Love and Care}
\label{ch:ka-love-and-care}
\hlchaptermodality{495}

\begin{chapteropening}
\textbf{By the end of this chapter:} \emph{I can say to cuddle, to kiss, tender affection, motherly love and care.}

Complete the last lesson of chapter 495: I can say to cuddle, to kiss, tender affection, motherly love and care.
\end{chapteropening}

\section[\texorpdfstring{\kn{ಮುದ್ದಿಸು} (muddisu)}{ಮುದ್ದಿಸು (muddisu)}]{\kn{ಮುದ್ದಿಸು} (muddisu) — to cuddle; to fondle}

\label{lesson:KA-C495-muddisu}



\begin{quote}
Before the new one: say the Kannada for to lose heart, then the Kannada for bewilderment.
\end{quote}

\subsection*{You'll want to know: \kn{ಮುದ್ದಿಸು}}
\textbf{\kn{ಮುದ್ದಿಸು}} — \emph{muddisu} — \textquotedblleft{}to cuddle; to fondle\textquotedblright{}.

\begin{quote}
\kn{ಅಜ್ಜಿ} \kn{ಮೊಮ್ಮಗನನ್ನು} \kn{ಮುದ್ದಿಸಿದರು}.
\end{quote}

\begin{quote}
\emph{ajji mommaganannu muddisidaru.}
\end{quote}

\begin{quote}
Grandmother cuddled her grandson.
\end{quote}

Past (I …): \kn{ಮುದ್ದಿಸಿದೆ}, \emph{muddiside}.

The noun \kn{ಮುದ್ದು} is fondness, and a pet word for a little one: \kn{ಮುದ್ದು} \kn{ಮಗು}, a darling child.

\begin{cousinweb}
\textbf{\kn{ಮುದ್ದಿಸು}} (\emph{muddisu}) is native Dravidian: DEDR 4960 lists Kannada \emph{muttu} and \emph{muddu}, \textquotedblleft{}a kiss\textquotedblright{}, with words for caressing and fondling; -\emph{isu} makes the verb.
\end{cousinweb}

\begin{grammarlens}[title={What you've built}]
One more word for describing what you lived and felt.
\end{grammarlens}

\subsection*{Your turn}
\begin{itemize}
\raggedright
  \item \emph{Say it:} \emph{muddisu}
  \item \emph{Say it:} \emph{muddisu}, once more
  \item \emph{Say it:} \emph{ajji mommaganannu muddisidaru.} — Grandmother cuddled her grandson.
  \item \emph{Recall:} say the Kannada for to lose heart, then the Kannada for bewilderment, then say \emph{muddisu} again
\end{itemize}

\begin{tcolorbox}[breakable,colback=teal!4,colframe=teal!35!black,title={Before you move on}]
What does \kn{ಮುದ್ದಿಸು} mean? (\textquotedblleft{}To cuddle; to fondle\textquotedblright{}.) Say it once more.
\end{tcolorbox}
\section[\texorpdfstring{\kn{ಮುತ್ತಿಡು} (muttiḍu)}{ಮುತ್ತಿಡು (muttiḍu)}]{\kn{ಮುತ್ತಿಡು} (muttiḍu) — to kiss}

\label{lesson:KA-C495-muttidu}



\begin{quote}
Before the new one: say the Kannada for bewilderment, then the Kannada for to cuddle.
\end{quote}

\subsection*{You'll want to know: \kn{ಮುತ್ತಿಡು}}
\textbf{\kn{ಮುತ್ತಿಡು}} — \emph{muttiḍu} — \textquotedblleft{}to kiss\textquotedblright{}.

\begin{quote}
\kn{ಅಮ್ಮ} \kn{ಮಗುವಿನ} \kn{ಕೆನ್ನೆಗೆ} \kn{ಮುತ್ತಿಟ್ಟಳು}.
\end{quote}

\begin{quote}
\emph{amma maguvina kennege muttiṭṭaḷu.}
\end{quote}

\begin{quote}
The mother kissed the baby's cheek.
\end{quote}

Past (I …): \kn{ಮುತ್ತಿಟ್ಟೆ}, \emph{muttiṭṭe}.

\kn{ಮುತ್ತು}, a kiss + \kn{ಇಡು}, to put: what is kissed takes -\kn{ಗೆ}. \kn{ಮುತ್ತು} is also a pearl.

\begin{cousinweb}
\textbf{\kn{ಮುತ್ತಿಡು}} (\emph{muttiḍu}) is a native Dravidian compound: \emph{muttu}, \textquotedblleft{}a kiss\textquotedblright{} (DEDR 4960), and \emph{iḍu}, \textquotedblleft{}to put, to place\textquotedblright{} (DEDR 442). DEDR also records \emph{muddiḍu}, \textquotedblleft{}to kiss\textquotedblright{}.
\end{cousinweb}

\begin{grammarlens}[title={What you've built}]
One more word for describing what you lived and felt.
\end{grammarlens}

\subsection*{Your turn}
\begin{itemize}
\raggedright
  \item \emph{Say it:} \emph{muttiḍu}
  \item \emph{Say it:} \emph{muttiḍu}, once more
  \item \emph{Say it:} \emph{amma maguvina kennege muttiṭṭaḷu.} — The mother kissed the baby's cheek.
  \item \emph{Recall:} say the Kannada for bewilderment, then the Kannada for to cuddle, then say \emph{muttiḍu} again
\end{itemize}

\begin{tcolorbox}[breakable,colback=teal!4,colframe=teal!35!black,title={Before you move on}]
What does \kn{ಮುತ್ತಿಡು} mean? (\textquotedblleft{}To kiss\textquotedblright{}.) Say it once more.
\end{tcolorbox}
\section[\texorpdfstring{\kn{ವಾತ್ಸಲ್ಯ} (vātsalya)}{ವಾತ್ಸಲ್ಯ (vātsalya)}]{\kn{ವಾತ್ಸಲ್ಯ} (vātsalya) — tender affection}

\label{lesson:KA-C495-vatsalya}



\begin{quote}
Before the new one: say the Kannada for to cuddle, then the Kannada for to kiss.
\end{quote}

\subsection*{You'll want to know: \kn{ವಾತ್ಸಲ್ಯ}}
\textbf{\kn{ವಾತ್ಸಲ್ಯ}} — \emph{vātsalya} — \textquotedblleft{}tender affection\textquotedblright{}.

\begin{quote}
\kn{ಅಜ್ಜನ} \kn{ವಾತ್ಸಲ್ಯ} \kn{ನನಗೆ} \kn{ಇನ್ನೂ} \kn{ನೆನಪಿದೆ}.
\end{quote}

\begin{quote}
\emph{ajjana vātsalya nanage innū nenapide.}
\end{quote}

\begin{quote}
I still remember grandfather's tender love.
\end{quote}

The love of an elder for a little one; Sanskrit \emph{vatsa} is a calf.

\begin{cousinweb}
\textbf{\kn{ವಾತ್ಸಲ್ಯ}} (\emph{vātsalya}) is borrowed from Sanskrit \emph{vātsalya}, \textquotedblleft{}affection\textquotedblright{} (CDIAL 11506), built on \emph{vatsala}, \textquotedblleft{}fond\textquotedblright{}, from \emph{vatsa}, \textquotedblleft{}a calf\textquotedblright{}.
\end{cousinweb}

\begin{grammarlens}[title={What you've built}]
One more word for describing what you lived and felt.
\end{grammarlens}

\subsection*{Your turn}
\begin{itemize}
\raggedright
  \item \emph{Say it:} \emph{vātsalya}
  \item \emph{Say it:} \emph{vātsalya}, once more
  \item \emph{Say it:} \emph{ajjana vātsalya nanage innū nenapide.} — I still remember grandfather's tender love.
  \item \emph{Recall:} say the Kannada for to cuddle, then the Kannada for to kiss, then say \emph{vātsalya} again
\end{itemize}

\begin{tcolorbox}[breakable,colback=teal!4,colframe=teal!35!black,title={Before you move on}]
What does \kn{ವಾತ್ಸಲ್ಯ} mean? (\textquotedblleft{}Tender affection\textquotedblright{}.) Say it once more.
\end{tcolorbox}
\section[\texorpdfstring{\kn{ಮಮತೆ} (mamate)}{ಮಮತೆ (mamate)}]{\kn{ಮಮತೆ} (mamate) — motherly love; deep attachment}

\label{lesson:KA-C495-mamate}



\begin{quote}
Before the new one: say the Kannada for to kiss, then the Kannada for tender affection.
\end{quote}

\subsection*{You'll want to know: \kn{ಮಮತೆ}}
\textbf{\kn{ಮಮತೆ}} — \emph{mamate} — \textquotedblleft{}motherly love; deep attachment\textquotedblright{}.

\begin{quote}
\kn{ಅಮ್ಮನ} \kn{ಮಮತೆಗೆ} \kn{ಬೆಲೆ} \kn{ಕಟ್ಟಲು} \kn{ಆಗುವುದಿಲ್ಲ}.
\end{quote}

\begin{quote}
\emph{ammana mamatege bele kaṭṭalu āguvudilla.}
\end{quote}

\begin{quote}
A mother's love cannot be priced.
\end{quote}

From Sanskrit \kn{ಮಮ}, mine: the feeling that someone is your very own.

\begin{cousinweb}
\textbf{\kn{ಮಮತೆ}} (\emph{mamate}) is borrowed from Sanskrit \emph{mamatā}, \textquotedblleft{}the state of mine, a sense of ownership, attachment\textquotedblright{} (Monier-Williams), built on \emph{mama}, \textquotedblleft{}mine\textquotedblright{}.
\end{cousinweb}

\begin{grammarlens}[title={What you've built}]
One more word for describing what you lived and felt.
\end{grammarlens}

\subsection*{Your turn}
\begin{itemize}
\raggedright
  \item \emph{Say it:} \emph{mamate}
  \item \emph{Say it:} \emph{mamate}, once more
  \item \emph{Say it:} \emph{ammana mamatege bele kaṭṭalu āguvudilla.} — A mother's love cannot be priced.
  \item \emph{Recall:} say the Kannada for to kiss, then the Kannada for tender affection, then say \emph{mamate} again
\end{itemize}

\begin{tcolorbox}[breakable,colback=teal!4,colframe=teal!35!black,title={Before you move on}]
What does \kn{ಮಮತೆ} mean? (\textquotedblleft{}Motherly love; deep attachment\textquotedblright{}.) Say it once more.
\end{tcolorbox}
\section[\texorpdfstring{\kn{ಕಾಳಜಿ} (kāḷaji)}{ಕಾಳಜಿ (kāḷaji)}]{\kn{ಕಾಳಜಿ} (kāḷaji) — care; concern}

\label{lesson:KA-C495-kalaji}



\begin{quote}
Before the new one: say the Kannada for tender affection, then the Kannada for motherly love.
\end{quote}

\subsection*{You'll want to know: \kn{ಕಾಳಜಿ}}
\textbf{\kn{ಕಾಳಜಿ}} — \emph{kāḷaji} — \textquotedblleft{}care; concern\textquotedblright{}.

\begin{quote}
\kn{ದಾದಿಗೆ} \kn{ರೋಗಿಗಳ} \kn{ಬಗ್ಗೆ} \kn{ತುಂಬಾ} \kn{ಕಾಳಜಿ} \kn{ಇದೆ}.
\end{quote}

\begin{quote}
\emph{dādige rōgigaḷa bagge tumbā kāḷaji ide.}
\end{quote}

\begin{quote}
The nurse cares a great deal about the patients.
\end{quote}

Care for someone: … \kn{ಬಗ್ಗೆ} \kn{ಕಾಳಜಿ}.

\begin{cousinweb}
\textbf{\kn{ಕಾಳಜಿ}} (\emph{kāḷaji}) is usually traced to Marathi \emph{kāḷjī}, \textquotedblleft{}care, anxious concern\textquotedblright{}; its earlier history is not settled.
\end{cousinweb}

\begin{grammarlens}[title={What you've built}]
One more word for describing what you lived and felt.
\end{grammarlens}

\subsection*{Your turn}
\begin{itemize}
\raggedright
  \item \emph{Say it:} \emph{kāḷaji}
  \item \emph{Say it:} \emph{kāḷaji}, once more
  \item \emph{Say it:} \emph{dādige rōgigaḷa bagge tumbā kāḷaji ide.} — The nurse cares a great deal about the patients.
  \item \emph{Recall:} say the Kannada for tender affection, then the Kannada for motherly love, then say \emph{kāḷaji} again
\end{itemize}

\begin{tcolorbox}[breakable,colback=teal!4,colframe=teal!35!black,title={Before you move on}]
What does \kn{ಕಾಳಜಿ} mean? (\textquotedblleft{}Care; concern\textquotedblright{}.) Say it once more.
\end{tcolorbox}
